\section{Scalar reference paths and admissible perturbations}
\label{sec:paths}

We construct the deterministic clocks used in the pressure comparison.
Their first requirement is that the reference field \(h_*=K-bq\)
be nonnegative and increasing. Since \(q(v)=\Gamma(K(v))\), its
increment along scalar time \(t=K(v)\) is
\[
 dh_*=(1-b\Gamma'(t))\,dt-q\,db.
\]
The construction first creates a positive margin in
\(1-b\Gamma'\), then chooses the matrix variation and the additional
field terms within that margin.

We use direct paths while varying the perturbation strength, and
smoothed paths while varying the size coordinate. The smoothing keeps
the terminal mass fixed and permits size differentiation; a final
comparison shows that its pressure cost is smaller than \(x^5\).
Put \(\ell=\exp(\theta/6)\),
\(0\le\theta\le\log2\). Direct paths have terminal mass \(a=\ell\),
and smoothed paths have terminal mass \(a=4\). The parameter \(\ell\)
is fixed except during a size change. In normalized scalar coordinates,
time is \(\xi=\ell^2t\) and mass is \(v/\ell\). Fix
\[
 \omega=x^{5+1/4},\quad M=\log^3(1/x),\quad
 x^{12}\le\lambda\le x^{p_*},\qquad
 U=\min(1,x\lambda^{-s_*/2}),\qquad \lambda_1=\lambda^{1+\eta}.
                                                               \tag{P1}\label{eq:P1}
\]
The final strength means \(\lambda=x^{p_*}\). Choose \(s_*<1/2\)
sufficiently close to \(1/2\), then \(\eta>0\) sufficiently small
also relative to \(1/2-s_*\), and \(p_*>0\) sufficiently small.
Later small exponents may be decreased subject to their stated order.
Logarithms and constants are absorbed with powers much smaller than
\(p_*\) times the available gaps between exponents of \(\lambda\).

\subsection{The fixed-temperature scalar reference}

For a one-site clock \(K\) of terminal mass 1, write
\(f_{\rm spin}(K)=f_1(0,K)\) for its self-corrected pressure. Let
\(z:(0,1)\to[0,1]\) be the deterministic nondecreasing quantile
minimizing the fixed-temperature scalar objective
\[
 f_{\rm spin}(Tz)+\frac T4\int_0^1z(w)^2\,dw.
\]
The corresponding variance-time quantile is \(Tz(w)\). Its distribution
function is
\(A_0(\xi)=|\{w\in(0,1):Tz(w)\le\xi\}|\), and \(\Gamma_0\)
is the scalar response of this clock as defined in
Section~\ref{sec:foundations}. These are the fixed base objects;
the reference clocks below will perturb them.

\begin{lemma}[Regular scalar profile and a strict top plateau]
\label{paths:base-profile}
At fixed \(T>1\), the scalar minimizing distribution has support \([0,\xi_*]\), satisfies \(\Gamma_0(\xi)=\xi/T\) there, and has a positive initial density. Its active-side distribution function has bounded one-sided derivatives of every fixed order. Moreover \(p_{\rm top}=A_0(\xi_*-)\) belongs to \((0,1)\).
\end{lemma}

\begin{proof}
We use PE, Lemma~5.8, whose full-support input is \cite[Theorem~1.1 and Lemma~2.2]{Lopatto2026}. Its measure has support \([0,\xi_*]\), \(\xi_*>0\); on this support \(\Gamma_0(\xi)=\xi/T\). Its hypotheses are precisely fixed \(T>1\), zero field, and the Ising
two-spin scalar objective above, with terminal mass 1. The lemma gives \(A_0(\xi)\sim \kappa_T\xi/T\) at 0, \(\kappa_T>0\), and continuity of the quantile on the open mass interval. The measure therefore has support down to 0 with no interior gap or atom at 0.

At the minimizing \(z\), vary \(z\) to \(z+tP_0(z)\) with smooth \(P_0\) supported in \((0,1)\) (both signs admissible for small \(t\)). By scalar \eqref{eq:A1} the derivative of the objective is \(T\int (z-\Gamma_0(Tz))P_0(z)/2=0\), giving self-consistency by continuity throughout the support. On \([0,\xi_*]\) viewed from the active side, \(A_0\) is smooth with bounded one-sided derivatives of all orders, \(A'_0\ge c>0\) near zero. Indeed \(\Gamma_0''=0\) in the interior gives
\(A_0=\mathbf E (V_{zzz})^2/(2\mathbf E(V_{zz})^3)\) almost everywhere by PE, Section~5, Lemma~5.3 (in variance time).

The numerator and denominator are continuous moment functions. The
positive-Hessian argument of the same lemma, restricted to a bounded
field interval of uniformly positive probability, bounds the denominator
away from zero throughout the compact active interval. Their ratio therefore
gives a continuous representative of \(A_0\), extending Lipschitzly
to both endpoints. Differentiation with the scalar PDE and It\^o
generator improves this representative one derivative at a time;
the positive denominator allows the iteration to every fixed order.
The positive initial slope is the initial-density conclusion of PE.

In fact \(p_{\rm top}=A_0(\xi_*-)\) belongs to \((0,1)\). Here is a useful argument. Extend by a mass-one annealed clock, and let \(\rho\) denote the law of \(Z_{\xi_*}\). The finite measure \(w(dz)=\rho(dz)/\cosh z\) is even and log-concave.
To see the sign of the potential, take a step approximation ending at
\(\xi_*\) and telescope its transition weights. After division by
\(\cosh z\), the joint density of the Gaussian increments is their
Gaussian density times the exponential of negative convex continuation
potentials. Each coefficient is a nonnegative mass increment; the
unused terminal mass similarly multiplies \(-\log\cosh z\).
Marginalizing the increments at fixed terminal field preserves
log-concavity~\cite{Prekopa1973}. Fixed-clock approximation, as in PE,
Proposition~2.14, passes it to \(w\). The limiting field law has a
density because its diffusion has bounded drift over positive time. At subsequent time \(\xi_*+d\), the derivative \(\mathbf E\chi^2\), \(\chi=V_{zz}\), is
\[
 e^{-d/2}\int\operatorname{sech}^3 z\,(w*G_d)(dz)
 \le e^{-d/2}\Gamma'_0(\xi_*),
\]
with \(G_d\) the centered Gaussian kernel. The inequality follows by even unimodality (translate overlaps of upper level intervals). At \(d=0\) the two sides agree, so the inequality makes the right
second derivative strictly negative. The left second derivative is
zero by \(\Gamma_0(\xi)=\xi/T\). The spatial moments in
\(\Gamma''=\mathbf E(V_{zzz}^2-2A_0\chi^3)\) agree across the
endpoint; only the mass changes. Consequently
\[
 \Gamma_0''(\xi_*+)-\Gamma_0''(\xi_*-)
 =-2(1-p_{\rm top})\mathbf E\chi^3<0,
\]
which gives \(p_{\rm top}<1\). Its positivity follows from the
positive initial density and monotonicity of \(A_0\). Thus the base
quantile is constant on a nonempty terminal mass interval.

\end{proof}

\subsection{Creating a positive scalar gap}

The minimizing scalar reference is marginal: \(\Gamma'_0=1/T\)
on its active support, so in these normalized coordinates the choice
\(b=T\) leaves no margin for \(-q\,db\). We perturb its density
near zero to create a gap of order \(\lambda t^2\). The matrix
construction will add a second margin of order \(\lambda U^2\).

Choose \(z_0>0\) small fixed with \(A'_0\ge c>0\) up to \(2z_0\), and let \(P(\xi)\) be smooth nonnegative, bounded by \(\xi\), equal to \(\xi\) on \([0,z_0/2]\) and zero after \(z_0\). Adding \(\lambda P\) gives the distribution
\(A_\lambda=A_0+\lambda P\), which is admissible for sufficiently
small \(\lambda\). Its direct physical clock is
\(K(v)=t(v)=\ell^{-2}A_\lambda^{-1}(v/\ell)\); it is constant
on the top plateau \(v/\ell\in[p_{\rm top},1]\).

For the smoothed clock, fix the final strength. Spread the jump
\(1-p_{\rm top}\) linearly over \([\xi_*,\xi_*+\omega]\), add
\(\omega\xi\) to this mass profile for \(0\le\xi<M\), and at
\(M\) jump to total normalized mass \(4/\ell\). Denote the resulting
profile by \(\widehat A\). Below \(M\) it is independent of
\(\ell\), has no support gaps, and satisfies
\(\omega\le\widehat A'\le C/\omega\) almost everywhere.
The physical smoothed clock is
\(K(v)=\ell^{-2}\widehat A^{-1}(v/\ell)\). It may have a terminal
interval on which its normalized quantile equals \(M\). The positive
density below that interval permits the size differentiations used later.

For either reference, changing from normalized to physical coordinates gives
\[
 V(t,z)=\ell^{-1}\widehat V(\xi,\ell z),\qquad
 \Gamma(t)=\widehat\Gamma(\xi),
\]
where the normalized mass profile was just given, with terminal mass \(j=1\) or \(j=4/\ell\). Define \(b_\circ=1/V_{zz}(0,0)^2\). Since the starting field is zero,
\(\Gamma'(0)=V_{zz}(0,0)^2\), so this choice sets the initial gap
\(1-b_\circ\Gamma'(0)\) to zero. The lift below makes that gap
positive at every positive time.

\begin{lemma}[Quantitative scalar lift]
\label{paths:lift}
Uniformly (extending by a terminal annealed clock),
\[
 \begin{array}{ll}
 b_\circ\asymp1,\quad
 1-b_\circ\Gamma'(t)\ \ge c\lambda\min(t^2,1)
     & (t\ge0),\\
 1-b_\circ\Gamma'(t)\ \le C\lambda t^2
     & (0\le \xi\le \xi_*;\ \text{active side}),\\
 t(v)\asymp v,\quad q(v)\asymp v
     & \text{near }0,\quad q(v)\asymp v\text{ throughout}.
 \end{array}                                                       \tag{P2}\label{eq:P2}
\]
\end{lemma}

\Needspace{6\baselineskip}
\begin{proof}\leavevmode\par
\paragraph{The direct lift.} Work in normalized coordinates. Up to \(\xi_*\), positive-order derivatives of continuations change by \(O(\rho_1)\), \(\rho_1=\lambda\int P\). Indeed the PDE for their undifferentiated difference has linear drift with bounded smooth spatial coefficients on a fixed time interval and source bounded with derivatives by \(C\lambda P\); differentiate its diffusion representation (bounded drift derivatives). Coupling \(Z\)'s by common Brownian noise gives \(L^k\) change \(O_k(\rho_1\sqrt{\xi})\) for \(\xi\le z_0\), using oddness of the gradients, then \(O_k(\rho_1\sqrt{z_0})\) for \(\xi\in[z_0,\xi_*]\); continuation derivatives are unchanged on the latter interval. All constants here are uniform as \(z_0\) is decreased. Thus on \([0,z_0]\),
\[
 \widehat\Gamma''(\xi)=-2\lambda P(\xi)\mathbf E\widehat\chi^3+O(\xi\rho_1)
\]
by oddness of third derivatives; \(\mathbf E\widehat\chi^3\) stays strictly bounded below. Near zero, \(P(\xi)=\xi\), so integration gives a drop of order
\(\lambda\xi^2\) from \(\widehat\Gamma'(0)\). By time \(z_0\)
the accumulated drop is of order \(\rho_1\), while the integrated
error is \(O(\rho_1z_0^2)\). On \([z_0,\xi_*]\) the continuation
itself is unchanged, and
\(\widehat\Gamma'=1/T+O(\rho_1\sqrt{z_0})\). Choosing \(z_0\)
small first makes this error smaller than the accumulated drop. This
proves both the near-zero order and persistence of the gap. Beyond \(\xi_*\) use the annealed inequality just proved (valid also for \(A_\lambda\)). This argument differentiated once in \(\log\lambda\) gives \(O(\lambda\xi^2)\) for the derivative at fixed small \(\xi\) of the normalized gap (the linearized PDE and diffusion estimates are identical). Also \(\ell^2 b_\circ=T+O(\lambda)\).

\paragraph{Smoothing.} On any fixed compact time interval, continuing the direct PDE annealed to \(M\) with terminal mass \(j\) instead of 1 changes values, modulo spatial constants, by \(O(e^{-cM})\). Indeed at a fixed time above the support the expectation of \(\cosh^ {1/j}(j(z+B))\), over the remaining variance, lies between the sum of the two individual exponential contributions and the same sum times \(1-O(e^{-cM})\) (with the common factor \(2^{-1/j}\)), by subadditivity and restriction to disjoint half-lines for \(B\). Earlier nonlinear heat transitions preserve uniform additive comparison. Changing the normalized mass to the smoothed one now costs \(O(\omega M^2)\) in values, by PDE comparison and bounded gradients. Spatial derivatives on the compact interval thus change by \(O_\epsilon(\omega^{1-\epsilon})\), any fixed finite orders, via higher-order difference quotients and uniform spatial derivative bounds; diffusions couple there with the same power accuracy. Consequently the error in the active \(\widehat\Gamma''\) bound above is \(o(\lambda)\xi\) even at small \(\xi\) (use oddness, the added mass \(\omega\xi\), and Brownian coupling at small time with factor \(\sqrt{\xi}\)).

Beyond the support compare \(\widehat\Gamma'\) to the direct annealed case on a sufficiently long fixed interval. Thereafter use \(\widehat\chi\) between \(c(1-\widehat V_z^2)\) and \(C(1-\widehat V_z^2)\), by the scalar case of \eqref{eq:A1} since all remaining masses are bounded and bounded below. The martingale diffusion coefficient of \(\widehat V_z\) is \(\widehat\chi\); It\^o's formula for \(\sqrt{1-\widehat V_z^2}\) gives exponential decay of its expectation (localize inside \((-1,1)\) and use nonnegativity), so \(\mathbf E\widehat\chi^2\) eventually stays small. This proves \eqref{eq:P2} also with smoothing.

More explicitly, for the scalar conditional mean \(m_t\) and \(f(m)=\sqrt{1-m^2}\), localization gives
\[
 \frac{d}{dt}\mathbf E f(m_t)
 =-\frac12\mathbf E\frac{\chi_t^2}{(1-m_t^2)^{3/2}}
 \le -c\,\mathbf E f(m_t)
\]
once the remaining masses are bounded below. The initial value is at most one, so this bound is uniform in the starting field. It also applies conditionally at the beginning of the long interval and supplies the exponential terminal-datum bound used next.

\end{proof}

The lift provides the gap needed for admissibility. To differentiate
pressures later, we also need regularity in the path parameters.
There are two different requirements: near zero, derivatives must retain
the quadratic vanishing of the scalar gap; over all masses, their
absolute integrals must remain controlled despite the small smoothing
density \(\omega\).

\subsection{Regularity in strength and size}

\begin{lemma}[Parameter derivatives of the scalar references]
\label{paths:parameter-regularity}
Let the direct and smoothed scalar references have the parameters in
\eqref{eq:P1}, with \(1\le\ell\le2^{1/6}\); the smoothed reference
is at final strength. The following bounds hold uniformly for sufficiently
large \(N\). Every derivative order and compact time interval is fixed
before \(N\to\infty\).

For the direct reference, the normalized inverse near zero has uniformly
bounded one-sided derivatives of every fixed order. Their
\(\log\lambda\) derivatives are \(O(\lambda)\), and
\[
 \partial_{\log\lambda}q(v)=O(\lambda\min(v,1)).
\]
At fixed sufficiently small normalized time \(\xi\), the
\(\log\lambda\) derivative of
\(1-\widehat\Gamma'(\xi)/\widehat\Gamma'(0)\) is
\(O(\lambda\xi^2)\).

For the smoothed reference, write \(\widehat V_j\) and
\(\widehat\Gamma_j\) to display the terminal mass \(j=4/\ell\).
On each fixed compact scalar-time interval \([0,R]\), and for every
fixed spatial order \(r\ge1\), there are constants \(C_r,c_r>0\)
such that
\[
 \sup_{0\le\xi\le R,\ z\in\mathbb R}
 \left|\partial_j\partial_z^r\widehat V_j(\xi,z)\right|
 \le C_r e^{-c_rM}.
\]
For a fixed sufficiently small \(\xi_0>0\), this implies the more
precise cancellation
\[
 \left|\partial_j\bigl(
       \widehat\Gamma'_j(\xi)-\widehat\Gamma'_j(0)\bigr)\right|
 \le C e^{-cM}\xi^2,
 \qquad 0\le\xi\le\xi_0.                              \tag{P2a}\label{eq:P2a}
\]
For a fixed sufficiently small mass \(v_0>0\), the smoothed reference
also satisfies, for each fixed \(r\ge0\),
\[
 \sup_{0\le v\le v_0}
 \bigl(|\partial_v^r t|+|\partial_v^r q|
       +|\partial_\theta\partial_v^r t|
       +|\partial_\theta\partial_v^r q|\bigr)\le C_r,
\]
with one-sided derivatives at zero. The physical initial spatial
derivatives obey the rescaling bound, for each fixed \(r\ge1\),
\[
 \left|\partial_\theta\left(
     \ell^{1-r}\partial_z^r V(0,0)\right)\right|
 \le C_r e^{-c_rM}.
\]

Over the full mass interval \((0,4)\), the smoothed clocks \(t,q\)
are absolutely continuous in \(\theta\), with
\[
 \int_0^4\bigl(|\partial_\theta t(v)|
                  +|\partial_\theta q(v)|\bigr)\,dv
 \le C M^C.                                            \tag{P2b}\label{eq:P2b}
\]
The constants in this integrated bound contain no inverse power of
\(\omega\).
\end{lemma}

\begin{proof}
For the direct reference, \(A_0'(0)>0\) gives the inverse bounds.
The perturbation of its inverse is confined to the initial interval,
and differentiating its inverse relation gives the stated
\(O(\lambda)\) bounds. The linearized PDE and coupled diffusion in
the proof of Lemma~\ref{paths:lift} give the bound on
\(\partial_{\log\lambda}q\). That proof also differentiates the
scalar gap only after subtracting its value at zero, which preserves
the factor \(\xi^2\).

For size variations, the normalized smoothed profile below \(M\) is
independent of \(\ell\); only its terminal datum changes through
\(j=4/\ell\). The diffusion representation of
\(\partial_j\widehat V_j\) has terminal datum
\(\partial_j[j^{-1}\log\cosh(jz)]\). It differs from the constant
\(\log2/j^2\) by at most
\(C\sqrt{1-\tanh^2(jz)}\). The conditional martingale estimate in
the preceding proof therefore makes its expectation exponentially
small, uniformly in the starting field on a fixed compact time interval.
Higher spatial derivatives are uniformly bounded by the step power rules with one
bounded smooth terminal insertion and its bounded derivatives; the
remaining slots are bounded spin derivatives controlled by PE,
Corollary~2.10. If the uniform value error is \(\varepsilon\), global
finite-difference interpolation between orders zero and any fixed
\(k>r\) bounds the derivative of order \(r\) by
\(C_{r,k}\varepsilon^{1-r/k}\). This proves the asserted exponential
bound uniformly over all starting fields. The calculation differentiates
only the terminal datum while the mass profile below \(M\) stays fixed.

To obtain \eqref{eq:P2a}, differentiate the scalar second-derivative
identity at fixed normalized time. The spatial derivatives just bounded
are exponentially small, \(\widehat V_z\) and
\(\widehat V_{zzz}\) are odd, and the normalized mass is \(O(\xi)\)
near zero. Differentiating the Brownian coupling gives a field change
\(O_p(e^{-cM}\sqrt\xi)\). Hence
\(\partial_j\widehat\Gamma''_j(\xi)=O(e^{-cM}\xi)\), and
integration from zero proves \eqref{eq:P2a}.
Repeated PDE and generator differentiation bounds the response time
derivatives near zero, where the normalized mass profile is smooth with
bounded derivatives. Applying the same calculation to the equations
linearized in \(j\) bounds the mixed derivatives
\(\partial_j\partial_\xi^r\widehat\Gamma_j\) there. Since
\(A_0'(0)>0\), composition with the
normalized inverse and the exact rescalings
\(\xi=\ell^2t\), \(\widehat A(\xi)=v/\ell\) gives the stated
one-sided mass derivatives and their size derivatives. Finally,
\(\ell^{1-r}\partial_z^r V(0,0)=\partial_z^r\widehat V_j(0,0)\),
and \(\partial_\theta j=-j/6\). Thus treating the normalized initial
spatial derivatives as constants during a size differentiation incurs
only exponentially small errors.

It remains to control these derivatives after integration over all
masses. Off the terminal plateau, the inverse relation is
\(\widehat A(\xi)=v/\ell\), with \(\xi=\ell^2t(v)\). Since
\(\partial_\theta\ell=\ell/6\), differentiation at fixed \(v\) gives
\[
 \partial_\theta t(v)
 =-\frac{t(v)}3-
   \frac{v}{6\ell^3\widehat A'(\ell^2t(v))}.
\]
The second term can be large pointwise, but under
\(dv=\ell\widehat A'(\xi)\,d\xi\) its inverse density cancels.
Both resulting integrals are \(O(M)\). On the terminal plateau,
\(t=M/\ell^2\) and \(\partial_\theta t=-t/3\). Its moving boundary
creates no jump term because the clock is continuous there. In
particular \(\int|\partial_\theta t(v)|\,dv\lesssim M^C\), with
no inverse power of \(\omega\).

For \(q(v)=\widehat\Gamma_j(\xi)\), there is also dependence on the
terminal mass \(j=4/\ell\). We bound it at fixed normalized time
\(\xi\) by representing \(\widehat\Gamma_j(\xi)\) as the expectation
of the product of two terminal spins whose branches separate at that
time. This test is bounded by 1 and does not depend on \(j\). In a
step approximation, differentiate the logarithm of its positive tree
density, telescoped as in PE, Proposition~2.14. Below \(M\) all
normalized masses are fixed. The derivative
\(\partial_j[j^{-1}\log\cosh(jz)]\) is uniformly bounded, and the
diffusion representation gives the same bound for interior continuation
derivatives. Their coefficients have total variation bounded by the
number of observed leaves times the terminal mass, independently of
the number of steps. The remaining terminal terms are bounded by
\(C(1+|Z_M^1|+|Z_M^2|)\). Hence the logarithmic density derivative
has \(L^1\) norm at most \(C(1+M)\), because bounded scalar drift
gives \(\mathbf E|Z_M^i|\lesssim1+M\). It follows that
\[
 |\partial_j\widehat\Gamma_j(\xi)|\le C(1+M),
 \qquad 0\le\xi\le M,
\]
uniformly in the step resolution and in \(\omega\). The resulting
uniform Lipschitz bound passes to continuous clocks by PE,
Proposition~2.14; it supplies absolute continuity and the same derivative
bound almost everywhere without requiring convergence of derivatives.

Off the plateau, the chain rule now reads
\[
 \partial_\theta q(v)
 =-\frac{v\,\widehat\Gamma_j'(\xi)}
          {6\ell\widehat A'(\xi)}
   -\frac j6\partial_j\widehat\Gamma_j(\xi).
\]
After integration in \(v\), the first term again loses its inverse
density. Its integral is bounded by
\(C\int_0^M\widehat\Gamma_j'(\xi)\,d\xi\le C\), since
\(\widehat\Gamma_j\) is increasing and lies in \([0,1]\). The
second term has integral \(O(1+M)\). On the plateau only this second
term remains. Thus \(q\), like \(t\), has an absolutely integrable
size derivative of norm \(O(M^C)\).

\end{proof}

\subsection{Constructing admissible matrix and field clocks}

We now use the scalar gap to allow a varying matrix clock. The positive
term from the lift pays for its increase, while a negative constant
offset supplies a further margin at the smallest masses. We then add
field increments, checking separately the single ramp that can decrease.

Choose a small fixed mass \(v_b>0\), such that inverse densities are regular with bounded derivatives and bounded above and below up to \(8v_b\). Take a smooth nondecreasing \(G(v)\) equal to zero on \([0,x]\), \(v^2\) on \([8x,v_b]\) and constant positive above \(2v_b\), \(0\le G'\le C v\). Use dyadic masses \(u=2^{-h}\) in \([16x,v_b/16]\), a fixed smooth nonnegative \(\phi\) supported on \([1/2,4]\), positive on \([1,2]\), and parameters \(\gamma=(\gamma_0,\gamma_u)\). These cutoffs are common across path segments on the dyad and independent of \(\ell\). Take
\[
 b(v)=b_\circ+c_b\lambda\left(G(v)-U^2-c_1\gamma_0(G(v)+U^2)
                     -c_1\sum_u\gamma_u u^2\phi(v/u)\right),              \tag{P3}\label{eq:P3}
\]
with sufficiently small fixed \(c_1,c_b>0\). The parameters \(\gamma\) may range in \([-4,4]\). They are held fixed
when taking a positive-law expectation; later we average their coordinates
independently over \([1/3,2/3]\). By bounded overlap, \(0\le b'(v)\le Cc_b\lambda v\), \(b\asymp1\), \(c_g=-\partial_{\gamma_0}b\asymp\lambda(v^2+U^2)\) throughout this bounded mass interval. Also \(h_*=K-bq\) is nonnegative nondecreasing, with
\[
 dh_*\ \ge c\lambda(v^2+U^2)\,dt,\qquad
 h_*(v)\lesssim\lambda(v^3+vU^2)\quad\text{for small }v .              \tag{P4}\label{eq:P4}
\]
Indeed, along the inverse clock, including constant pieces,
\[
 dh_*=(1-b\Gamma'(t))\,dt-q\,db.
\]
The lift \eqref{eq:P2} supplies the \(\lambda v^2\) gap near zero,
and the negative \(U^2\) offset in \eqref{eq:P3} supplies
\(\lambda U^2\). The term \(q\,db\) occurs only at small active
mass, where \(q\asymp v\) and the inverse density is regular; choosing
\(c_b\) small absorbs it in the positive gap. Outside that interval
\(db=0\). At large \(t\), \(\Gamma'\) is small and the gap is
bounded below by a constant. This proves the first bound in
\eqref{eq:P4}; integrating its small-time upper counterpart from zero
gives the second. We also have for every \(z\ge0\) (with either weight \(w(y)=1\) or \(1-y\))
\[
 \int_0^1 w(y)\,[1-b(v)\Gamma'(t(v)+y b(v)(z-q(v)))]\,dy
           \ge c\lambda(v^2+U^2).                                        \tag{P5}\label{eq:P5}
\]
Indeed the first half of the secant lies at time at least \(t(v)/2\) since \(h_*\ge0\), hence the base lift there dominates \(c\lambda v^2\). The subtracted \(U^2\) offset in \eqref{eq:P3} provides the \(U^2\) term unless \(\Gamma'\) itself is small; everywhere the loss relative to these bounds from \eqref{eq:P3} is at most \(C c_b\lambda v^2\).

Set \(h(v)=0\) for \(v<x\); above that cut use
\[
\begin{gathered}
 h=h_*+H_{\rm lo}+\chi_{\rm lin}H_{\rm lin}+\chi_e H_e,\\
 H_{\rm lo}=\lambda_1U^3\log\frac{\min(v,\ell)}x,\qquad
 H_{\rm lin}=x^3(\min(v,\ell)-4v_b)_+.
\end{gathered}                                                        \tag{P6}\label{eq:P6}
\]
Here \(H_{\rm lo}\) supplies additional field increase above the lower
cut. The multipliers \(\chi_{\rm lin},\chi_e\) lie in \([0,1]\).
The linear term is introduced only at final strength and remains during
the size comparison. The term \(H_e\) is used only on direct paths;
it supplies field increase along their top plateau and is removed
before switching to smoothed paths. To define it, we join an active-side
ramp to a profile on that plateau. Take a sufficiently small fixed \(U_0>0\), \(\widetilde U=\min(U_0,U)\), a smooth taper \(\Upsilon(U)\) equal to 1 on \([0,U_0/2]\), zero on \([U_0,1]\), with values in \([0,1]\). Define \(E_b>0\) by
\[
 p_{\rm top}-A_0(\xi_*-E_b)=\widetilde U^4/E_b^2=d_-,\qquad d_+=\widetilde U^4.
\]
It exists uniquely near zero by strict increase, with \(|dE_b/d\log\lambda|\lesssim E_b\). On the plateau \(w=v/\ell\in[p_{\rm top},1]\) set
\[
 H_e=\lambda_1\Upsilon(U)\,\widetilde U^2
       \left((1-w+d_+)^{-1/2}-(w-p_{\rm top}+d_-)^{-1/2}\right).
                                                                        \tag{P7}\label{eq:P7}
\]
On the active part ramp linearly in \(\xi=\ell^2 t(v)\) from zero at \(\xi_*-E_b\) to its value at the plateau; below the ramp set \(H_e=0\). Thus on each dyadic strip of distance \(\asymp d'\) from an edge of the plateau up to its midpoint (using also one capped strip of width \(\asymp d_\pm\) adjacent to each edge), the clock increase per unit mass due to \(H_e\), for small \(U\), is comparable to \(\lambda_1 U^2/(d')^{3/2}\). On the active-side ramp its derivative in \(t\) may be negative,
but its magnitude is \(O(\lambda_1)\). This ramp lies at a fixed
positive mass, where \eqref{eq:P4} gives an increase at least
\(c\lambda\,dt\). Since \(\lambda_1/\lambda=\lambda^\eta\),
the negative part is absorbed for large \(N\). Thus the total field
clock remains nonnegative and nondecreasing even where the added defect
is negative; all paths are admissible. Both \(H_e\) and its \(\log\lambda\) derivative are bounded, on the ramp by \(C\lambda_1 E_b\), on the plateau by \(C\lambda_1\widetilde U^2/\sqrt{d'}\), vanishing away from \(U\le U_0\). In differentiating, \(\ell\) is then fixed and the inverse reference near the ramp does not vary.

\subsection{Defect bounds, endpoint tests, and parameter monotonicity}

In particular the field defect \(d=h-h_*\) satisfies
\[
 \int d^2/c_g\,dv\ \lesssim_\circ\
   \lambda^{1+2\eta}U^5+\chi_e^2\lambda^{1+2\eta}U^4+\chi_{\rm lin}^2 x^6/\lambda.       \tag{P8}\label{eq:P8}
\]
To verify \eqref{eq:P8}, separate the part below the cut from the
three added field terms. Below \(x\), \(d=-h_*\), and
\eqref{eq:P4} gives
\[
 \int_0^x\frac{d^2}{c_g}\,dv
 \lesssim\lambda U^2x^3
 \lesssim\lambda^{1+2\eta}U^5.
\]
The last inequality uses \(x/U\lesssim\lambda^{1/12}\) and sufficiently
small \(\eta\). For \(H_{\rm lo}\), integration of
\((v^2+U^2)^{-1}\) contributes \(O(U^{-1})\); after squaring its
amplitude, the bound is \(\lambda^{1+2\eta}U^5\) with logarithmic
loss. On the active ramp for \(H_e\), the squared amplitude is at
most \(C\lambda_1^2E_b^2\), its mass length is \(\ell d_-\), and
\(c_g\asymp\lambda\). The identity
\(E_b^2d_-=\widetilde U^4\) therefore gives
\(O(\lambda^{1+2\eta}U^4)\). On the plateau, the squares of the two
reciprocal-square-root terms integrate with only logarithmic loss and
give the same bound. Finally, \(H_{\rm lin}\) is supported at mass
bounded away from zero, where \(c_g\gtrsim\lambda\), so its contribution
is \(O(x^6/\lambda)\). The inequality for the square of a sum combines
these contributions with the displayed multipliers.

\paragraph{Strength derivatives.}
During a direct strength sweep, \(\ell,\gamma\) are fixed,
\((\chi_{\rm lin},\chi_e)=(0,1)\), and a dot means
\(\partial_{\log\lambda}\). The differentiated defects obey the
same localized bounds. On the mask they are
\(O(\lambda(v^3+vU^2))\): differentiate the scalar gap only after
subtracting its value at zero. Above the mask, differentiating
\eqref{eq:P6}--\eqref{eq:P7} gives the claimed ramp and plateau bounds.
In this regime all dotted clocks, \(\dot q\), and \(\dot b\) are
bounded by \(O_\circ(\lambda)\).

\paragraph{Size derivatives.}
On the smoothed size segment, the strength is final,
\((\chi_{\rm lin},\chi_e)=(1,0)\), and a dot means
\(\partial_\theta\), with \(\gamma\) fixed. Exact scalar rescaling
and Lemma~\ref{paths:parameter-regularity} give the same small-mass
defect bounds for \(d\) and \(\dot d\). In particular,
\eqref{eq:P2a} retains the factor \(\xi^2\) after differentiating
the scalar gap relative to its time-zero value. Above small
mass, \(d,\dot d=O_\circ(x^3)\). The other dotted clocks are
absolutely integrable with \(O_\circ(1)\) norms by the global
parameter bounds proved above.

\paragraph{Direct comparison endpoints.}
Rescale a direct path to \([0,1]\) by \(w=v/\ell\). The PE comparison
hypotheses at every lower scale \(w\gtrsim x\), with
\(T_m=\ell^2b_\circ\to T\), take the form
\[
 \ell^2 b(\ell w)=T_m+o(w^2),\qquad \ell^2 h(\ell w)=o(w^3),\qquad w\gtrsim x,
                                                                        \tag{P9}\label{eq:P9}
\]
because \(\lambda U^j\le x^j\lambda^{1-js_*/2}\) for \(j=2,3\).
The gaps \(1-js_*/2>0\) make these errors vanish relative to the
required powers at the lower comparison scales.

In varying \(\gamma\) it is useful to subtract \(\int b q^2/4\) from \(f_m^0\). Its increasing derivative along any matrix decrease \(c=-\partial_\gamma b\ge0\) accompanied by the path prescription satisfies
\[
 \partial_\gamma(f_m^0-\int bq^2/4)
 \ \ge\ \tfrac14\int c\,(D_v^2+(S_v-q)^2)\,dv,                         \tag{P10}\label{eq:P10}
\]
To check the sign, note that \(q\) is fixed under the \(\gamma\)
variation. Above the mask the prescribed field satisfies
\(\partial_\gamma h=cq\), so \eqref{eq:A1} gives the integrand
\[
 \frac c4(B_v-2qS_v+q^2)
 =\frac c4\big(D_v^2+(S_v-q)^2\big).
\]
On the mask, \(h=0\) remains fixed and the actual integrand is
\(c(B_v+q^2)/4\). It exceeds the displayed square by
\(cqS_v/2\ge0\), proving \eqref{eq:P10}.

\subsection{Interchanging the direct and smoothed paths}

\begin{proposition}[Comparison of the two regularizations]
\label{paths:switch}
At final strength \(\lambda=x^{p_*}\), with
\((\chi_{\rm lin},\chi_e)=(1,0)\), the direct and smoothed systems
satisfy \eqref{eq:P11}, uniformly with the same parameters
\(\gamma,\ell\). Their scalar trial gaps
\[
 f_{\rm spin}(K)+\frac14\int b^{-1}(K-h)^2\,dv-f_m^0,
\]
where \(f_{\rm spin}\) is the self-corrected one-site pressure at
the corresponding terminal mass, also differ by \(O(x^{5+c})\).
The same fixed \(c>0\) works for these three comparisons, independently
of the site count.
\end{proposition}

\begin{proof}
Fix \(\lambda=x^{p_*}\), \(\chi_{\rm lin}=1\), and \(\chi_e=0\),
and keep \(\gamma,\ell\) the same in both systems. First extend the
direct system to length 4 by a single field jump at mass \(\ell\).
For \(v>\ell\), put \(K=K_{\rm post}=M/\ell^2,\ q=1,\ b=b_{\rm post}=b(\ell-),\ h=h_{\rm post}=K_{\rm post}-b_{\rm post}+H_{\rm lo}(\ell)+H_{\rm lin}(\ell)\). On this extension \(q=1\) is a deterministic comparison function,
rather than a claim about the finite-time scalar response. We first
compare these four extended functions to their smoothed counterparts.
For fixed \(0<\epsilon<1/21\), put
\[
 E_\epsilon=\omega M^2+\omega^{1-\epsilon}+e^{-c'M}.
\]
Then
\[
 \begin{split}
 &\|K_{\rm hard}-K_{\rm smoothed}\|_1
 +\|q_{\rm hard}-q_{\rm smoothed}\|_1\\
 &\qquad+\|b_{\rm hard}-b_{\rm smoothed}\|_1
 +\|h_{\rm hard}-h_{\rm smoothed}\|_1
 \lesssim E_\epsilon.
 \end{split}
\]
Here and below the norms integrate over \(v\in(0,4)\).
For the inverse clocks,
the area between monotone quantiles equals the area between their
mass profiles; the spread jump and density floor contribute
\(O(\omega M^2)\). Below mass \(\ell\), both scalar times remain
in a fixed compact interval, so the earlier
\(O_\epsilon(\omega^{1-\epsilon})\) moment and initial-derivative
comparisons control the response at a common scalar time and
\(b_\circ\). The remaining change in \(q\), caused by evaluating
the response at different scalar times, is bounded by
\(\sup|\Gamma'|\,|K_{\rm hard}-K_{\rm smoothed}|\). Its integral
is therefore controlled by the inverse-clock estimate.

Above \(\ell\), the portion before the smoothed terminal plateau has
mass length \(O(\omega M)\). Outside that strip, the martingale bound
gives \(q=1-O(e^{-c'M})\). Equations \eqref{eq:P3} and \eqref{eq:P6}
then give the asserted \(L^1\) estimates for \(b,h\) as well.
Since \(\omega=x^{5+1/4}\), the choice \(\epsilon<1/21\) gives
\((5+1/4)(1-\epsilon)>5\). With \(M=\log^3(1/x)\), this proves
\(E_\epsilon=O(x^{5+c})\) for a fixed \(c>0\).

We next compute the pressure of this extension, denoted by
\(f_{\rm hard}\). Condition on arbitrary earlier Hamiltonian values
\((H_\sigma)_{\sigma\in\{-1,1\}^m}\), and set
\(D=h_{\rm post}-h(\ell-)\), \(p=\ell/4\). If
\(z\) has independent standard Gaussian coordinates and \(\Phi\)
is the standard Gaussian distribution function, then
\[
 \begin{split}
 &\Phi(\ell\sqrt D)^m e^{m\ell^2D/2}
          \sum_\sigma e^{\ell H_\sigma}\\
 &\qquad\le
 \mathbf E_z\left(\sum_\sigma
       e^{4H_\sigma+4\sqrt D\,z\cdot\sigma}\right)^p
 \le e^{m\ell^2D/2}\sum_\sigma e^{\ell H_\sigma}.
 \end{split}
\]
The upper bound is subadditivity for \(0<p\le1\). For the lower
bound, retain the term indexed by \(\sigma\) on the orthant
\(\sigma_i z_i\ge0\) for every \(i\). These orthants are disjoint,
and exponential tilting of each Gaussian coordinate gives the factor
\(\Phi(\ell\sqrt D)\). Thus the bound is uniform in the conditioned
Hamiltonian. After logarithms and division by \(m\ell\), the error
lies between \(\ell^{-1}\log\Phi(\ell\sqrt D)\) and zero. Here
\(D=M/\ell^2-O(1)\), so the error is \(O(e^{-c'M})\), independently
of \(m\).

Restoring the normalized spin prior multiplies the quantity inside this
expectation by \(2^{-mp}\), whereas the direct partition uses
\(2^{-m}\). Their contribution to the per-site difference is
\(\log2(1/\ell-1/4)\). The Gaussian expectation contributes
\(\ell D/2\). Combining this with the direct and length-four self
corrections, and propagating the resulting uniform additive comparison
through the earlier continuations, gives
\[
 f_{\rm hard}= f_{\rm direct}
      +(\ell-4)(b_{\rm post}/4+h_{\rm post}/2)
      +\log2\,(1/\ell-1/4)+O(e^{-c'M}).
\]
The calculation also holds for the tilted initial revelation, so its
explicit constant cancels in \(\psi\). For the corrected untilted
pressure, subtract the added tail \((4-\ell)b_{\rm post}/4\) of
\(\int bq^2/4\). The non-logarithmic correction becomes
\[
 \frac{\ell-4}{2}(b_{\rm post}+h_{\rm post})
 =\frac{\ell-4}{2}
   (K_{\rm post}+H_{\rm lo}(\ell)+H_{\rm lin}(\ell)),
\]
which is independent of \(\gamma\), as is the \(\log2\) term.
Every earlier power mean preserves a uniform additive comparison,
including the tilted initial continuation. Finally, \eqref{eq:A1}
bounds a pressure change by the \(L^1\) change of its clocks.
Applying it to the hard and smoothed systems yields,
uniformly,
\[
\begin{split}
 \psi_{\rm smoothed}&=\psi_{\rm direct}+O(x^{5+c}),\\
 (f_m^0-\int bq^2/4)_{\rm smoothed}
   &=(f_m^0-\int bq^2/4)_{\rm direct}
       +\mathrm{const}_{\gamma\text{-indep.}}+O(x^{5+c}).
\end{split}                                                          \tag{P11}\label{eq:P11}
\]
For the scalar trial gap we also need stability of its quadratic
integral. Although the smoothed \(K,h\) can be of order \(M\),
their difference is uniformly bounded: above the mask,
\(K-h=bq-H_{\rm lo}-H_{\rm lin}\), while below it
\(K-h=K=O(x)\). The hard extension has the same bound.
Since \(b\asymp1\), the function \((K-h)^2/b\) is uniformly
Lipschitz on these paths. Thus its integral changes by
\(O(E_\epsilon)\) under smoothing, just as the scalar and matrix
pressures do.

It remains to compute the hard-extension change. The scalar trial gap
\(f_{\rm spin}(K)+\int b^{-1}(K-h)^2/4-f_m^0\), where \(f_{\rm spin}\) is the self-corrected one-site pressure, changes by
\[
 (4-\ell)b_{\rm post}\big((K_{\rm post}-h_{\rm post})/b_{\rm post}-1\big)^2/4
       +O(x^{5+c})=O(x^{5+c})
\]
by the same hard-jump computation applied to the scalar pressure.
Indeed \(K_{\rm post}-h_{\rm post}=b_{\rm post}
-H_{\rm lo}(\ell)-H_{\rm lin}(\ell)\), so the displayed term is
\[
 \frac{4-\ell}{4b_{\rm post}}
       (H_{\rm lo}(\ell)+H_{\rm lin}(\ell))^2=O(x^6).
\]
At final strength,
\[
 H_{\rm lo}(\ell)
 =O\!\left(x^3\lambda^{1+\eta-3s_*/2}\log(1/x)\right)=o(x^3),
\]
and
\(H_{\rm lin}(\ell)=O(x^3)\). Decreasing the fixed saving \(c\)
if necessary incorporates this term in \(O(x^{5+c})\). The parameters,
in particular \(\ell\), have been kept identical throughout.

\end{proof}
